\subsection{Addition und Subtraktion}
\label{subsec:vector_algebra_addition_subtraction}

\begin{DefinitionBox}{Vektoraddition und Vektorsubtraktion}{vector_algebra_addition_subtraction}
  Für $\vec a,\vec b\in\R^n$ werden Addition und Subtraktion
  komponentenweise ausgeführt:
  \[
    \vec a+\vec b
    =\begin{pmatrix}a_1\\\vdots\\a_n\end{pmatrix}
     +\begin{pmatrix}b_1\\\vdots\\b_n\end{pmatrix}
    =\begin{pmatrix}a_1+b_1\\\vdots\\a_n+b_n\end{pmatrix},
  \]
  \[
    \vec a-\vec b
    =\vec a+(-\vec b)
    =\begin{pmatrix}a_1-b_1\\\vdots\\a_n-b_n\end{pmatrix}.
  \]
\end{DefinitionBox}

\begin{center}
  \CachedTikzFigure{content/resources/chapters/vector_algebra/tikz/vector_addition_subtraction.tex}
\end{center}

\begin{TheoremBox}{Rechenregeln der Vektoraddition}{vector_algebra_addition_laws}
  Für $\vec a,\vec b,\vec c\in\R^n$ gilt
  \[
    \begin{aligned}
      \vec a+\vec b&=\vec b+\vec a
        &&\text{(Kommutativgesetz)},\\[0.5em]
      \vec a+(\vec b+\vec c)&=(\vec a+\vec b)+\vec c
        &&\text{(Assoziativgesetz)},\\[0.5em]
      \vec a+\vec 0&=\vec a,
      &\qquad \vec a+(-\vec a)&=\vec 0.
    \end{aligned}
  \]
\end{TheoremBox}

\begin{ExampleBox}{Vektoren komponentenweise verrechnen}{vector_algebra_componentwise_operations}
  Gegeben seien
  \[
    \vec a=\begin{pmatrix}2\\0\\4\end{pmatrix},
    \qquad
    \vec b=\begin{pmatrix}-1\\2\\-10\end{pmatrix}.
  \]
  Dann gilt
  \[
    \vec a+\vec b
    =\begin{pmatrix}2+(-1)\\0+2\\4+(-10)\end{pmatrix}
    =\Cloze{3cm}{\begin{pmatrix}1\\2\\-6\end{pmatrix}},
    \qquad
    \vec a-\vec b
    =\begin{pmatrix}2-(-1)\\0-2\\4-(-10)\end{pmatrix}
    =\Cloze{3cm}{\begin{pmatrix}3\\-2\\14\end{pmatrix}}.
  \]
\end{ExampleBox}
